% revision/old-draft-notes.tex -------------------------------------------
% Salvage record for the two dissolved early drafts: paper/main.tex (deleted
% in eb2d896) and revised_manuscript/ (created in 41144dd, deleted in
% eb2d896). Plain notes file: NOT \input by revision/main.tex, never compiled.
%
% Recovered with:
%   git show HEAD~7:paper/main.tex
%   git show HEAD~7:revised_manuscript/main.tex
%   git show HEAD~7:revised_manuscript/chapters/*.tex
% (line numbers below are those of the HEAD~7 blobs; the content is unchanged
% from commit 41144dd, where the adaptation draft was created).
%
% What this file holds (paper/README.md, "Removed 2026-10-04"):
%   - the longer abstract variant of paper/main.tex (nothing else preserves
%     it: the validated submission keeps only the shorter variant);
%   - the nine passages that are genuinely unique to the adaptation draft,
%     checked sentence by sentence against paper/submission-metro/submission.tex
%     ("the baseline for what is novel", paper/README.md);
%   - the two title suggestions of the draft era.
%
% Where the rest of the dissolved material lives (nothing else is lost):
%   - six of the nine passages below are ALSO embedded as
%     "CARRIED OVER FROM AN OLD DRAFT" comments at their point of use in
%     revision/sections/(2),(4),(5),(6) (see revision/README.md,
%     "Carried-over material"); they are repeated here so the salvage is
%     complete in one place;
%   - the same CARRIED-OVER blocks also quote three passages that are NOT
%     draft-unique -- the housing-case empirical passage in (2), the
%     demand-side/supply-side view cluster in (4) and the labor-slack
%     conclusion in (6) are validated-submission text -- so they are not
%     reproduced here;
%   - the title suggestions also sit as a comment above \papertitle in
%     revision/main.tex;
%   - the figures of revised_manuscript/pictures/ went to paper/pictures/
%     (+ archive/), the bibliography is byte-identical to revision/references.bib.


%% Longer abstract

% source: paper/main.tex:83 (git show HEAD~7:paper/main.tex)
% The longer of the two abstract variants of the pluralist submission. The
% shorter variant that was actually submitted is preserved in
% paper/submission-metro/submission.tex:85; the two differ only in this one
% paragraph. (The adaptation draft's own, completely rewritten abstract is
% preserved below as novel passages 7-9.)

This paper revisits paradigmatic differences in economics with a focus on input-output modeling. We show that such differences represent deep divides in economic theorizing and impact strongly on the magnitude and signs of expected effects. At the same time estimates from input-output models are crucial for understanding the expected economic impacts of additional investment undertaken in the course of a socio-ecological transformation aiming to render social provisioning processes carbon-neutral. Taking the transformation of the German housing sector as a practical example this paper illustrates the divergence between typical results obtained and explores two central axiomatic variations of a neoclassical approach -- the introduction of CES-functions as well as labor slack -- that promise to explore some middle ground in between established approaches. Thereby we hope to better illuminate how different axiomatic setups imprint on estimates of the economic effects of ecologically motivated transformation efforts to provide applied researchers with better guidance when it comes to choosing foundational model assumptions.


%% Novel passages (9)

% (1) source: revised_manuscript/chapters/ch03_cge_model.tex:115 (git show
%     HEAD~7:revised_manuscript/chapters/ch03_cge_model.tex). Draft-unique
%     span only: the sentence opening that line ("Eventually, the question
%     emerges how to link such a setup with the Input-Output tables ...") is
%     validated-submission text and is not repeated. Also embedded at
%     revision/sections/(4)kernel-closures-financings.tex:42 (CARRIED OVER).

In practice, such a link necessitates that $\Omega$ plays a role in the exact specification of the production function and/or consumption function. For a Cobb-Douglas function, as the most widely used specification, the cost shares contained in $\Omega$ are interpreted as output elasticities and, correspondingly, used to specify the relative contribution of factors and intermediates to net output of a given sector, which translates into

% (2) source: revised_manuscript/chapters/ch03_cge_model.tex:161-165. The
%     Laspeyres real-GDP measurement of the adaptation draft; the submission
%     (and paper/main.tex) use the Tornqvist/Divisia passage at
%     paper/submission-metro/submission.tex:419 instead, so the two are
%     alternatives, not copies. Also embedded at
%     revision/sections/(4)kernel-closures-financings.tex:44 (CARRIED OVER).

Calculating GDP directly using these post-shock prices would yield a nominal GDP measure. To obtain real GDP, one typically employs a price index such as the Laspeyres or Paasche index. We use the Laspeyres index, so that GDP is computed as
$$
\text{GDP} = \frac{\sum_{s=1}^n p_s \post c_s}{\sum_{s=1}^n p_s  c_s}
$$
which measures the change in consumption valued at pre-shock prices, providing a meaningful comparison of real economic activity before and after the shock.

% (3) source: revised_manuscript/chapters/ch05_ces_functions.tex:71.
%     Draft-unique in full. Also embedded at
%     revision/sections/(5)application.tex:42 (CARRIED OVER).

An overall inspection of the results show that in this comparison the standard Leontief approach delivers the largest GDP-estimates -- on this basis GDP is expected to grow by about 1.5\%, which is equivalent to a multiplier of 1.15 given that the initial investment impulse was specified to be 40.337 bn. €, about 1.3\% of GDP. In contrast, the Cobb-Douglas specification predicts a small reduction in GDP that results from efficiency-losses arising from the increase in green investment. These losses are due to the fact that the imposed green investment leads to a forced reallocation of intermediates and production factors towards the expanding sectors, which alters an already optimal allocation and, hence, induces inefficiencies. In turn, real output decreases slightly and the major macroeconomic impact of a green investment demand shock in this framework is to induce sector-specific inflationary pressures.

% (4) source: revised_manuscript/chapters/ch05_ces_functions.tex:100.
%     Draft-unique span only: the two sentences opening that line ("Finally,
%     by assessing changes in sectoral production after the demand shock ...
%     which focuses on final demand only).") are validated-submission text.
%     Also embedded at revision/sections/(5)application.tex:43 (CARRIED OVER).

This reproduces the finding from section \ref{subsec:CGE_elasticity} that a demand shock will depress real GDP on a sectoral level as the (comparatively lower) average expansion of shocked sectors does under no conditions fully compensate for the reduction in the production of unaffected sectors. Notwithstanding this overall pattern, the complex substitution dynamics implied the CGE approach lead to exceptional outcomes in specific sectors. The most obvious idiosyncratic outcome occurs in two sectors that are not shocked –– \textit{building construction works (sector 33 )} and \textit{civil engineering works (sector 34)} --, but, rather, massively expand their intermediate goods production at the expense of final outputs, which represents an over-proportional intersectoral crowding-out effect. In other words, intermediary goods from these sectors are in extremely high demand due to the enforced shift in final expenditures. For both of these sectors the expansion in intermediate goods production is so strong, that these sectors expand more strongly than in the Leontief case and thereby partially compensate for the lack of freely available workers in shocked sectors.

% (5) source: revised_manuscript/chapters/ch05_ces_functions.tex:73.
%     Draft-unique in full. Also embedded at
%     revision/sections/(6)beyond-demand-only.tex:44 (CARRIED OVER).

With regard to how changes in assumptions on substitution elasticities impact on final outcomes Figure \ref{fig:elasticity-gradient} indicates that real GDP is typically increasing in all three elasticities employed -- the higher the chosen elasticities, the higher are the respective GDP-estimates. This finding reproduces the intuition that easier substitution will lead to higher (aggregate) output and, hence, CES-based estimates are consistently lower than their Cobb-Douglas counterparts.

% (6) source: revised_manuscript/chapters/ch05_ces_functions.tex:75.
%     Draft-unique span only: the sentence opening that line ("Against this
%     backdrop we concede that the introduction of Leontief-like properties in
%     the production function helps bridging paradigmatic cleavages on a level
%     of conceptual foundations -- e.g. with respect to the question how to
%     adequately represent production processes.") is validated-submission
%     text. Also embedded at revision/sections/(6)beyond-demand-only.tex:46
%     (CARRIED OVER).

However, we find that introducing such a conceptual shift in neoclassical IO-models will actually serve to increase the gap between both approaches in terms of predicted outcomes. This result is not too surprising when taking into account that the assumption of a Leontief technology implies less flexibility as compared to a Cobb-Douglas technology. Hence, in the context of assumptions on substitutability neoclassical approaches are more optimistic than the heterodox archetype: viewed in isolation greater substitutability comes with higher prospects for growth in case of demand shocks and, correspondingly, leads to a more positive assessment of fiscal intervention in terms of forced green investment. However, this relatively more optimistic stance does in no way compensate for the fact that the estimated growth effects remain negative, which is effectively driven by assumptions on capacity constraints. However, before turning to this key aspect in section \ref{sec:labor} further below, we first inspect structural differences between model outcomes on a more granular level.

% (7) source: revised_manuscript/main.tex:81 (the adaptation draft's rewritten
%     abstract, paragraph 1). Not embedded anywhere in revision/sections/;
%     the only copy.

It is well understood that the IO multiplier model is a special case of a CGE model with fixed prices and unconstrained factor supply \citep{Robinson.2006}. The labor supply elasticity~$\eta$ --- running from zero (full-employment CGE) to infinity (IO multiplier) --- is thus the key parameter governing the relationship between these modeling traditions. While this conceptual bridge has been recognized for decades \citep{Rose.1995,Dervis.1982}, its quantitative implications have not been systematically measured.

% (8) source: revised_manuscript/main.tex:83 (the adaptation draft's rewritten
%     abstract, paragraph 2). Not embedded anywhere in revision/sections/;
%     the only copy.

We replicate and extend the Baqaee--Farhi framework to quantify the relative importance of factor market closures versus production technology parameters. Using the German housing transformation toward carbon neutrality as a policy-relevant application, we conduct a systematic sensitivity analysis across the labor supply elasticity continuum and a range of CES substitution elasticities. We find that the factor market closure explains roughly an order of magnitude more variation in both aggregate and sectoral outcomes than all production technology parameters combined. The estimated GDP effect of the housing transformation ranges from virtually zero to approximately 1.5\% depending on the closure, while CES elasticities shift this estimate by only a fraction of a percentage point. Moving from exogenous labor slack to endogenous labor supply produces the ``bridge'' the literature already knows exists --- and here we show how much it matters.

% (9) source: revised_manuscript/main.tex:85-87 (front matter of the same
%     abstract; counted with it as one passage unit). Not embedded anywhere
%     in revision/sections/; the only copy.

\textbf{Keywords:} input-output modeling, general equilibrium, labor supply elasticity, factor market closures, CES production functions, socio-ecological transformation

\textbf{JEL-Codes:} C67, C68, D57, E24


%% Title suggestions

% The two title suggestions of the draft era --- the pair that
% revision/main.tex:84-87 records above \papertitle and that paper/README.md
% refers to as "both title suggestions". Only (2) originates in the deleted
% revised_manuscript/; (1) is the working title of paper/narrative_outline.md
% (live, lines 46-47), paired with (2) when the revision was set up.

% (1) source: paper/narrative_outline.md:46-47 (the outline's working title);
%     quoted as suggestion (1) in revision/main.tex:85.

Which Short Run? Which Money? Closure Choice in Input-Output Models of the Socio-Ecological Transformation

% (2) source: revised_manuscript/main.tex:68-69 (git show
%     HEAD~7:revised_manuscript/main.tex); devised for the adaptation draft
%     (the subtitle idea goes back to the retired salvage plan,
%     docs/archive/REVISED_SALVAGE_PLAN.md:163). Quoted as suggestion (2) in
%     revision/main.tex:85-86.

Labor supply elasticities and sectoral dynamics in the input-output modeling continuum: a quantitative decomposition using the German housing transformation

% For completeness: a third title string of the same early-draft line sits in
% paper/superseded/abstract.typ:8-10 (frozen, kept visible per ADR-0004):
% "Substitution, Slack, and Scarcity: Navigating Paradigmatic Divides in
% Input-Output Analysis for Socio-ecological Transformation".
